\chapter*{Résumé}
\thispagestyle{empty}

Ce projet de fin d'études porte sur la conception, le développement,
le déploiement et l'analyse de sécurité d'une application web de prise
de notes personnelle, baptisée \textbf{WebCyber}. L'application repose
sur le \emph{framework} Python \textbf{Flask} pour la couche métier,
sur \textbf{PostgreSQL} pour la persistance des données, et sur
\textbf{Nginx} comme proxy inverse assurant la terminaison TLS.
L'ensemble est \textbf{conteneurisé} avec Docker et orchestré via Docker
Compose, puis \textbf{déployé sur AWS} (instance EC2 \texttt{t2.micro}
sous Ubuntu Server 22.04 LTS), avec une résolution de nom de domaine
gérée par \textbf{Route\,53} et un certificat HTTPS automatique fourni
par \textbf{Let's Encrypt}.

\vspace{0.4em}

Une attention particulière a été portée à la sécurité, organisée
en \textbf{six couches} : pare-feu réseau (\emph{Security Group} AWS),
durcissement du système, isolation Docker, configuration TLS et
en-têtes de sécurité Nginx, contrôles applicatifs (authentification,
hachage \texttt{PBKDF2-SHA256}, protection CSRF, validation des
entrées, isolation des données par utilisateur) et sécurité de la
base de données. La solution déployée a ensuite été soumise à une
campagne d'évaluation à l'aide de \textbf{Nmap}, \textbf{Nikto} et
\textbf{testssl.sh}, dont les résultats sont analysés et discutés.

\vspace{0.4em}

\textbf{Mots-clés :} application web, sécurité, cybersécurité, AWS,
EC2, Docker, conteneurs, Nginx, Flask, PostgreSQL, HTTPS, TLS,
Let's Encrypt, Nmap, Nikto, défense en profondeur, DevOps.

\vspace{1.5em}

\begin{center}
\noindent\rule{0.6\linewidth}{0.4pt}
\end{center}

\vspace{0.8em}

\section*{Abstract}

This end-of-studies project covers the design, development, deployment
and security assessment of a personal note-taking web application named
\textbf{WebCyber}. The application is built on the Python
\textbf{Flask} framework for business logic, \textbf{PostgreSQL} for
data persistence, and \textbf{Nginx} as a reverse proxy handling TLS
termination. The whole stack is \textbf{containerised} with Docker and
orchestrated through Docker Compose, then \textbf{deployed on AWS}
(EC2 \texttt{t2.micro} instance running Ubuntu Server 22.04 LTS), with
DNS resolution provided by \textbf{Route\,53} and HTTPS certificates
automatically issued by \textbf{Let's Encrypt}.

\vspace{0.4em}

Security has been treated as a first-class concern and structured into
\textbf{six layers}: AWS Security Group network firewall, host
hardening, Docker isolation, Nginx TLS and security headers,
application-level controls (authentication, \texttt{PBKDF2-SHA256}
password hashing, CSRF protection, input validation, per-user data
isolation) and database security. The deployed solution was then
assessed using \textbf{Nmap}, \textbf{Nikto} and \textbf{testssl.sh},
and the results are analysed and discussed.

\vspace{0.4em}

\textbf{Keywords:} web application, security, cybersecurity, AWS, EC2,
Docker, containers, Nginx, Flask, PostgreSQL, HTTPS, TLS, Let's Encrypt,
Nmap, Nikto, defense-in-depth, DevOps.
